\documentclass[10pt,a4paper,ssfamily]{exam}
\usepackage[brazil]{babel}
\usepackage[numbers,sort&compress]{natbib}
\usepackage[utf8]{inputenc}
\renewcommand{\baselinestretch}{1.0}
\usepackage{epsfig}
\usepackage{mhchem}
\usepackage{graphicx}
\usepackage{svg}
\usepackage{makeidx}
\usepackage{multicol}
\usepackage{multirow}
\usepackage{amssymb}
\usepackage{amsmath}
\DeclareMathOperator{\sen}{sen}
\newcommand{\listfont}{\bf\fontencoding{OT1}\sffamily\large}
\newcommand{\titlesfont}{\fontencoding{OT1}\sffamily\large}
\renewcommand{\baselinestretch}{1.0}
\renewcommand{\t}{\textrm}
\newcommand{\cc}[1]{[\textrm{#1}]}
\usepackage{enumitem}
\setlist{nosep, topsep=2pt, partopsep=0pt, parsep=0pt, itemsep=2pt}
\newcommand{\nomera}{
  \noindent
  \begin{tabular}{|p{0.7\textwidth}|p{0.3\textwidth}|}
  Nome: & RA: \\
  \hline
  \end{tabular}
}

\newcommand{\qini}{\begin{enumerate}[resume]\item}
\newcommand{\qend}{\end{enumerate}}

\newcommand{\oC}{$^\circ$C}
\newcommand{\e}[1]{$\times 10^{#1}$}

\begin{document}
\sffamily
\pagestyle{empty}
\nomera

\begin{center}
{\bf QG101 - Química Geral}\\
{\bf Bloco 13 - Termoquímica Básica}\\
\underline{Justifique suas respostas.}
\end{center}

\vspace{0.2cm}
\begin{center}{\bf Questões}\end{center}

% ---------------------------------------------------------------
\qini
Classifique cada processo abaixo como \textbf{exotérmico} ou
\textbf{endotérmico}, justificando em uma frase:
\begin{enumerate}
  \item Combustão de madeira.
  \item Derretimento de gelo.
  \item Fotossíntese.
  \item Respiração celular.
\end{enumerate}

\vspace{0.3cm}
\qend

% ---------------------------------------------------------------
\qini
Uma reação química libera $50$~kJ para a vizinhança.
\begin{enumerate}
  \item Qual é o sinal de $\Delta H$ dessa reação?
  \item A reação é exotérmica ou endotérmica? Justifique.
  \item Esboce, em palavras, como seria o diagrama de energia
        (reagentes x produtos) dessa reação.
\end{enumerate}

\vspace{0.3cm}
\qend

% ---------------------------------------------------------------
\qini
Calcule a quantidade de calor necessária para aquecer $100$~g de água de
$25$\oC{} a $75$\oC{} (calor específico da água
$= 4{,}18$~J~g$^{-1}$\oC$^{-1}$). Expresse o resultado em joules e em
quilojoules.

\vspace{0.8cm}
\qend

% ---------------------------------------------------------------
\qini
Dada a equação termoquímica:
$$\ce{C(s) + O2(g) -> CO2(g)} \qquad \Delta H = -394~\text{kJ/mol}$$
\begin{enumerate}
  \item O que significa o sinal negativo de $\Delta H$?
  \item Quantos kJ são liberados na combustão completa de $2{,}00$~mol de
        \ce{C(s)}?
  \item Quantos kJ são liberados na combustão completa de $6{,}0$~g de
        \ce{C(s)} ($M = 12{,}0$~g/mol)?
\end{enumerate}

\vspace{0.3cm}
\qend

% ---------------------------------------------------------------
\qini
Por que a transpiração resfria o corpo? Explique em termos de processo
endotérmico/exotérmico e do tipo de mudança de estado envolvida.

\vspace{0.8cm}
\qend

% ---------------------------------------------------------------
\qini
São conhecidos os seguintes dados termoquímicos:
\begin{align*}
\ce{H2(g) + 1/2 O2(g) &-> H2O(l)} &\Delta H_1 &= -286~\text{kJ/mol}\\
\ce{H2(g) + 1/2 O2(g) &-> H2O(g)} &\Delta H_2 &= -242~\text{kJ/mol}
\end{align*}
\begin{enumerate}
  \item Usando a Lei de Hess, calcule o $\Delta H$ da vaporização da água,
        $\ce{H2O(l) -> H2O(g)}$.
  \item O valor obtido deveria ser positivo ou negativo? Esse resultado é
        coerente com o tipo de processo (fornecimento ou liberação de
        calor)?
\end{enumerate}

\vspace{0.3cm}
\qend

% ---------------------------------------------------------------
\qini
Uma amostra de $100$~g de água e uma amostra de $100$~g de ferro
($c = 0{,}45$~J~g$^{-1}$\oC$^{-1}$), inicialmente à mesma temperatura,
recebem a mesma quantidade de calor ($5{,}0$~kJ cada).
\begin{enumerate}
  \item Calcule a variação de temperatura ($\Delta T$) de cada amostra.
  \item Qual delas esquenta mais? Relacione o resultado ao conceito de
        calor específico e a um exemplo do cotidiano (ex.: areia da praia
        x água do mar).
\end{enumerate}

\vspace{0.3cm}
\qend

% ---------------------------------------------------------------
\qini
A síntese da amônia é descrita por:
$$\ce{N2(g) + 3H2(g) -> 2NH3(g)} \qquad \Delta H = -92~\text{kJ/mol}$$
\begin{enumerate}
  \item Qual o valor de $\Delta H$ para a formação de \textbf{1 mol} de
        \ce{NH3(g)} a partir de \ce{N2} e \ce{H2}?
  \item Escreva a equação termoquímica da reação \textbf{inversa}
        (decomposição de $2$~mol de \ce{NH3}) e indique seu $\Delta H$.
\end{enumerate}

\vspace{0.3cm}
\qend

% ---------------------------------------------------------------
\qini
Um gás, ao absorver $300$~J de calor da vizinhança, se expande e realiza
$120$~J de trabalho sobre a vizinhança.
\begin{enumerate}
  \item Identifique o sinal de $q$ e de $w$ nesse processo, usando a
        convenção $\Delta U = q + w$ (em que $w$ é o trabalho realizado
        \emph{sobre} o sistema).
  \item Calcule $\Delta U$ do gás.
\end{enumerate}

\vspace{0.3cm}
\qend

% ---------------------------------------------------------------
\qini
\textbf{Calorimetria.} Um pedaço de $50{,}0$~g de um metal, inicialmente a
$100{,}0$\oC, é colocado em $100{,}0$~g de água a $20{,}0$\oC. O sistema
atinge o equilíbrio térmico a $24{,}0$\oC. Use
$c_{\text{água}} = 4{,}18$~J~g$^{-1}$\oC$^{-1}$ e admita que todo o calor
perdido pelo metal é ganho pela água (sistema isolado).
\begin{enumerate}
  \item Calcule o calor ganho pela água.
  \item Calcule o calor específico do metal, sabendo que ele perde a
        mesma quantidade de calor que a água ganha.
\end{enumerate}

\vspace{0.3cm}
\qend

\end{document}
